\LevelZero{Language}\label{ch:language}

\LevelOne{Definition}
The files inside \tbf{config/fuji/languages/} directory are named \tbf{language files}. \\
Each \tbf{language file} defines a table to map a \tbf{language key} into a \tbf{language value}.

\begin{tips}{List all languages.}
    Issue \cmd{/fuji inspect languages} command, to see what language files are loaded and in use.
\end{tips}

\LevelOne{Modify the default language.}
You can modify the \tbf{default language} in \mainconfig file. \\
Then, issue \cmd{/fuji reload} command to apply the changes.

\LevelOne{Control how the text is used.}
The \tbf{language value} only defines \str{what is the text}. \\
It doesn't define \str{how the text is used}. \\
You can insert \tbf{language instructions} to specify how to use that \tbf{language value}. \\

\begin{note}{Supported language instrutions}
    \begin{enumerate}
        \item {If the language value starts with \str{[send-action-bar]}, the text will be displayed in the action bar.}
        \item {If the language value starts with \str{[send-main-title]}, the text will be displayed in the main title.}
        \item {If the language value starts with \str{[send-sub-title]}, the text will be displayed in the sub title.}
        \item {If the language value starts with \str{[suppress-sending]}, the text will not be displayed.}
        \item {Otherwise, the text will be displayed in the chat message.}
    \end{enumerate}
\end{note}


\LevelOne{Control the text streaming}
For each fuji command, you can specify \str{--silent true/false} and \str{--stdout true/false} flags to control the text streaming during the command execution.\\

\begin{example}{}
    \begin{enumerate}
        \item {The \cmd{/fly others Steve --silent true} command will \tbf{not display the text} during the execution of \cmd{/fly} command.}
        \item {The \cmd{/fly others Steve --stdout true} command will \tbf{stream the text into the console} during the execution of \cmd{/fly} command.}
        \item {The \cmd{/run as console --stdout true send-message @r Hello} command will pass the flags into the \cmd{/send-message @r Hello} sub-command.}
    \end{enumerate}
\end{example}


